\documentclass[11pt,a4paper]{article}
\usepackage{xcolor}
\usepackage[utf8]{inputenc}
\usepackage{geometry}
\geometry{margin=1in}
\usepackage{hyperref}
\usepackage{enumitem}
\usepackage{listings}

\lstset{
  language=C,
  basicstyle=\ttfamily\small,
  keywordstyle=\bfseries,
  commentstyle=\itshape,
  showstringspaces=false,
  breaklines=true,
  captionpos=b,
  numbers=left,
  numberstyle=\tiny\color{gray},
}

\title{Digital Design Training Program \\ \large Day 9: Data Structures and Parsing Algorithms}
\author{Authored by }
\date{April 1, 2026}

\begin{document}

\maketitle

\section*{Theme: Translating human-readable equations into a machine-processable format using fundamental computer science algorithms and memory-safe C structures.}

\subsection*{Module 9.1: Embedded Memory \& Static Stacks (45 Minutes)}

\begin{itemize}
    \item \textbf{Goal:} Understand the dangers of dynamic memory allocation (\verb|malloc|/\verb|free|) on a microcontroller and build a memory-safe, statically allocated stack.
    \item \textbf{Conceptual Overview:} On a PC with gigabytes of RAM, asking the operating system for memory (\verb|malloc|) is common. On a microcontroller with just 2KB of SRAM, this is a recipe for disaster. The ``Heap'' (where \verb|malloc| gets memory) can become fragmented, leaving small, unusable gaps. Eventually, a request for memory will fail, causing a silent, hard-to-debug crash.
    \item \textbf{The Solution: Static Allocation.} We will pre-allocate all the memory our program will ever need when we compile it. For our calculator's stacks, we will use fixed-size arrays. This is memory-safe and predictable.
    \item \textbf{Action Steps \& Code:}
\begin{lstlisting}[caption={Code for Module 9.1: A Memory-Safe Static Stack}]
#include <stdio.h> // For NULL

#define STACK_SIZE 32 // Define a hard limit for our stack

// A C structure to hold our stack's data and state.
typedef struct {
    float data[STACK_SIZE]; // The fixed-size array for data
    int top;                // Index of the top element, -1 if empty
} FloatStack;

// Initialize the stack. Takes a POINTER to the stack struct.
void fstack_init(FloatStack *s) {
    s->top = -1;
}

// Push a value onto the stack.
// Returns 0 on success, -1 on failure (stack full).
int fstack_push(FloatStack *s, float value) {
    if (s->top >= STACK_SIZE - 1) {
        return -1; // Error: Stack Overflow
    }
    s->data[++s->top] = value;
    return 0;
}

// Pop a value from the stack.
// Returns the value, or 0.0f if the stack was empty.
float fstack_pop(FloatStack *s) {
    if (s->top < 0) {
        return 0.0f; // Error: Stack Underflow
    }
    return s->data[s->top--];
}

int main(void) {
    // Instantiate the stack. It lives on the main() call stack.
    FloatStack my_stack;
    
    // Pass a POINTER to the stack to our functions.
    // This is efficient; we don't copy the whole 32-element array.
    fstack_init(&my_stack);
    
    fstack_push(&my_stack, 3.14f);
    fstack_push(&my_stack, 1.618f);
    
    // float val = fstack_pop(&my_stack); // val is now 1.618f
    
    while(1) {}
    return 0;
}
\end{lstlisting}
    \item \textbf{Checkpoint:} You can create and use a custom, statically-allocated stack, passing it by reference (pointer) to functions to efficiently and safely manipulate its contents.
\end{itemize}
\newpage
\subsection*{Module 9.2: Lexical Analysis (Tokenization) (45 Minutes)}

\begin{itemize}
    \item \textbf{Goal:} Convert a raw string of characters from the keypad into a sequence of meaningful mathematical ``tokens''.
    \item \textbf{Conceptual Overview:} A computer doesn't understand ``12.5 + sin(3)''. It first needs to break this string into a list of ``words'' or tokens. This process is called lexical analysis or tokenization.
    \begin{center}
    \verb|"12.5+sin(3)"| $\rightarrow$ \verb|[NUMBER:"12.5"]|, \verb|[OPERATOR:"+"]|, \verb|[FUNCTION:"sin"]|, \verb|[LPAREN:"("]|, \verb|[NUMBER:"3"]|, \verb|[RPAREN:")"]|
    \end{center}
    \item \textbf{Action Steps \& Code:}
\begin{lstlisting}[caption={Code for Module 9.2: A Simple Tokenizer}]
#include <string.h>
#include <ctype.h> // For isdigit, isalpha, isspace

// An enum to define all possible token types
typedef enum {
    T_NONE, T_NUMBER, T_OPERATOR, T_FUNCTION, T_LPAREN, T_RPAREN
} TokenType;

// A struct to hold one token
typedef struct {
    TokenType type;
    char value[16]; // The string representation of the token
} Token;

// The tokenizer function
int tokenize(const char* input_str, Token* tokens, int max_tokens) {
    int count = 0;
    const char* p = input_str;

    while (*p != `\0' && count < max_tokens) {
        if (isspace(*p)) { p++; continue; } // Skip whitespace

        // Group digits and `.' into a NUMBER token
        if (isdigit(*p) || *p == `.') {
            const char* start = p;
            while (isdigit(*p) || *p == `.') p++;
            int len = p - start;
            tokens[count].type = T_NUMBER;
            strncpy(tokens[count].value, start, len);
            tokens[count].value[len] = `\0';
            count++;
            continue;
        }

        // Handle single-character operators
        if (strchr(``+-*/^'', *p)) {
            tokens[count].type = T_OPERATOR;
            tokens[count].value[0] = *p;
            tokens[count].value[1] = `\0';
            count++; p++;
            continue;
        }
        
        // ... Add similar blocks for parentheses and functions (isalpha) ...
        
        p++; // Failsafe for unknown characters
    }
    return count;
}
\end{lstlisting}
    \item \textbf{Checkpoint:} Given a mathematical expression as a C string, your \verb|tokenize| function can successfully produce an array of \verb|Token| structs, correctly identifying and classifying each part of the expression.
\end{itemize}
\newpage
\subsection*{Module 9.3: Dijkstra's Shunting-Yard Algorithm (90 Minutes)}

\begin{itemize}
    \item \textbf{Goal:} Convert the tokenized ``infix'' expression into ``Reverse Polish Notation'' (RPN) to easily solve the order of operations problem.
    \item \textbf{Conceptual Overview:} Humans write \verb|3 + 4 * 2|. A simple computer prefers \verb|3 4 2 * +|. This is RPN. It can be solved linearly with a single stack: read a token; if it's a number, push it to the stack. If it's an operator, pop two numbers, perform the operation, and push the result back. The Shunting-Yard algorithm is the classic method for converting from infix to RPN.
    \item \textbf{The Algorithm:} You need one queue for your RPN output and one stack for operators.
    \begin{enumerate}
        \item Read a token.
        \item If it's a number, push it to the output queue.
        \item If it's a function or a left parenthesis, push it to the operator stack.
        \item If it's an operator (e.g., \verb|+|): While the operator on top of the stack has greater or equal precedence, pop it to the output queue. Then, push the current operator to the stack.
        \item If it's a right parenthesis: Pop operators from the stack to the queue until you find a left parenthesis. Discard both parentheses. If the token on top of the stack is now a function, pop it to the queue.
        \item When all tokens are read, pop any remaining operators from the stack to the queue.
    \end{enumerate}
    \item \textbf{Action Steps \& Code:}
\begin{lstlisting}[caption={Code for Module 9.3: Shunting-Yard Implementation}]
// --- Assume Tokenizer and a Token-based Stack are defined ---

// Function to get operator precedence
int get_precedence(const Token* t) {
    if (t->type == T_FUNCTION) return 4;
    if (t->type != T_OPERATOR) return 0;
    switch (t->value[0]) {
        case `^': return 3;
        case `*': case `/': return 2;
        case `+': case `-': return 1;
        default: return 0;
    }
}

// The main shunting-yard function
int shunt(Token* infix, int infix_len, Token* rpn_queue) {
    int rpn_len = 0;
    TokenStack op_stack;
    tstack_init(&op_stack);

    for (int i = 0; i < infix_len; i++) {
        Token t = infix[i];
        if (t.type == T_NUMBER) {
            rpn_queue[rpn_len++] = t; // Rule 2
        } else if (t.type == T_FUNCTION || t.type == T_LPAREN) {
            tstack_push(&op_stack, t); // Rule 3
        } else if (t.type == T_OPERATOR) {
            // Rule 4
            while (op_stack.top > -1 && get_precedence(&tstack_peek(&op_stack)) >= get_precedence(&t)) {
                rpn_queue[rpn_len++] = tstack_pop(&op_stack);
            }
            tstack_push(&op_stack, t);
        } else if (t.type == T_RPAREN) {
            // Rule 5
            while (op_stack.top > -1 && tstack_peek(&op_stack).type != T_LPAREN) {
                rpn_queue[rpn_len++] = tstack_pop(&op_stack);
            }
            tstack_pop(&op_stack); // Discard left parenthesis
            if (op_stack.top > -1 && tstack_peek(&op_stack).type == T_FUNCTION) {
                rpn_queue[rpn_len++] = tstack_pop(&op_stack);
            }
        }
    }
    // Rule 6: Pop remaining operators
    while (op_stack.top > -1) {
        rpn_queue[rpn_len++] = tstack_pop(&op_stack);
    }
    return rpn_len;
}
\end{lstlisting}
    \item \textbf{Checkpoint:} Your \verb|shunt| function can take the token array for \verb|3+4*2/(1-5)^2| and correctly produce the RPN token array: \verb|3 4 2 * 1 5 - 2 ^ / +|.
\end{itemize}

\end{document}
